\section{Đóng góp của nghiên cứu}

Thông qua việc giải quyết các mục tiêu đặt ra, khóa luận mang lại ba đóng góp chính cả về mặt phương pháp luận khoa học lẫn giá trị ứng dụng thực tiễn trong miền y tế tiếng Việt:

Đóng góp thứ nhất nằm ở việc **Chuẩn hóa Khung hệ thống RAG cho Y khoa Tiếng Việt**. Nghiên cứu đề xuất, cài đặt và đánh giá thành công một đường ống RAG tiêu chuẩn (Standard Baseline RAG Pipeline) kết hợp giữa mô hình nhúng văn bản và mô hình ngôn ngữ lớn, được tối ưu hóa cho miền dữ liệu chuyên khoa Cơ Xương Khớp bằng tiếng Việt. Đường ống này thiết lập một cột mốc so sánh (Baseline) được thực nghiệm bài bản, đóng vai trò nền tảng khoa học cho việc đánh giá và so sánh hiệu quả của các kỹ thuật RAG nâng cao trong tương lai.

Đóng góp thứ hai là việc **Xây dựng Bộ dữ liệu và Bộ chuẩn đánh giá chuyên ngành (Vietnamese Orthopedic Benchmark)**. Nghiên cứu thực hiện quy trình thu thập, tiền xử lý và chuẩn hóa bộ ngữ liệu tri thức y khoa chuyên sâu về các bệnh lý cơ xương khớp từ các nguồn tài liệu y học chính thống tại Việt Nam. Đồng thời, khóa luận thiết lập một bộ câu hỏi đánh giá thực nghiệm (Clinical Evaluation Queries) được xây dựng theo định hướng các bộ chuẩn y tế quốc tế uy tín \cite{xiong2024mirage}, đóng góp một tài nguyên dữ liệu có giá trị cho cộng đồng nghiên cứu NLP trong y tế tại Việt Nam.

Đóng góp thứ ba là việc **Thực nghiệm và Đánh giá định lượng toàn diện**. Khóa luận cung cấp các bằng chứng thực nghiệm rõ ràng thông qua việc đo lường định lượng hệ thống dựa trên tập chỉ số RAG Metrics chuẩn hóa (bao gồm Faithfulness, Context Precision, và Context Recall) kết hợp với quy trình kiểm định chuyên môn bởi các chuyên gia y tế. Kết quả nghiên cứu chứng minh khả năng giảm thiểu đáng kể rủi ro ảo giác của RAG so với các mô hình ngôn ngữ lớn thuần túy, khẳng định tính hiệu quả và độ tin cậy của hệ thống trong tư vấn thông tin y học.
